\documentclass[11pt]{article}
\usepackage{amsmath,amssymb,booktabs,geometry,graphicx,hyperref,url}
\usepackage{subcaption,placeins}
\hypersetup{hidelinks}
\geometry{margin=1in}
\setlength{\parskip}{0.45em}
\setlength{\parindent}{0pt}
\graphicspath{{./}{../}}
\title{Stage 11: Integrated Modular Workflow With PSOR Audit}
\author{Free-Boundary-Aware Risk Surfaces for American Options}
\date{\today}
\begin{document}
\maketitle

\section{Stage 11: Integrated Modular Workflow With PSOR Audit}

\subsection{Purpose}
This stage evaluates the modular workflow formed by separate price, boundary,
and Delta heads. Scenario screening means producing rapid preliminary
estimates for a parameter regime and then auditing selected or unfamiliar cases
with CN/PSOR. The workflow is not a replacement for the benchmark and is not
evidence that one neural network solves price, the free boundary, and Greeks.

\subsection{Method}
The workflow loads three selected artifacts: the positive-premium price model, the direct boundary diagnostic head, and the bounded supervised Delta head. Six held-out test regimes are selected from the Stage 6 split table: three puts and three dividend-paying calls. For each regime, fresh CN/PSOR labels are regenerated on the Stage 6 grid before evaluating the workflow.

\subsection{Mathematical Formulation}
The price component predicts nonnegative continuation premium,
\[
\widehat{U}=\Phi+\widehat{P}_{\theta}, \qquad \widehat{P}_{\theta}\geq 0.
\]
The boundary component predicts a separate supervised estimate $\widehat{B}/K$. The Delta component predicts a bounded supervised estimate,
\[
\widehat{\Delta}_{\mathrm{put}}\in[-1,0], \qquad
\widehat{\Delta}_{\mathrm{call}}\in[0,1].
\]
These outputs are modular diagnostics evaluated against regenerated PSOR labels.

\subsection{Implementation Details}
The workflow loader is implemented in \path{src/workflow.py}. The integrated experiment is implemented in \path{experiments/13_integrated_application_workflow.py}. The experiment writes case metrics, by-regime metrics, violation tables, timing tables, fresh PSOR cases, and figures under the Stage 11 result directories.

\subsection{Validation and Diagnostics}
\begin{table}[h]
\centering
\begin{tabular}{llrrr}
\toprule
Case & Type & Price RMSE & Boundary MAE & Strict Delta RMSE\\
\midrule
put\_T0p25\_sig0p2\_r0p1\_q0p03 & Put & 0.0081 & 0.1100 & 0.0461\\
put\_T0p25\_sig0p4\_r0p01\_q0 & Put & 0.0108 & 0.1455 & 0.0231\\
put\_T0p25\_sig0p4\_r0p01\_q0p06 & Put & 0.0114 & 0.3270 & 0.0230\\
call\_T0p25\_sig0p2\_r0p05\_q0p1 & Call & 0.0058 & 0.1229 & 0.0443\\
call\_T0p25\_sig0p4\_r0p05\_q0p06 & Call & 0.0095 & 0.1296 & 0.0218\\
call\_T0p25\_sig0p4\_r0p1\_q0p1 & Call & 0.0096 & 0.0884 & 0.0216\\
\bottomrule
\end{tabular}
\caption{Price, boundary, and strict-mask Delta errors for six deterministic held-out regimes regenerated with CN/PSOR. The small audit set supports a workflow demonstration, not a broad generalization claim.}
\end{table}

\begin{table}[h]
\centering
\begin{tabular}{lr}
\toprule
Metric & Value\\
\midrule
Mean price RMSE & 0.0092\\
Mean boundary MAE & 0.1539\\
Mean strict Delta RMSE & 0.0300\\
Maximum obstacle violation & 0.0000\\
Maximum Delta-bound violation & 0.0000\\
Mean speedup & 87.52\\
\bottomrule
\end{tabular}
\caption{Aggregate metrics over the six fresh PSOR audit regimes. Speedup is PSOR label-generation time divided by neural workflow evaluation time and should be interpreted together with the accuracy columns.}
\end{table}

\subsection{Results}
The integrated workflow preserves the hard obstacle and Delta-bound checks on the evaluated cases. Price and strict-mask Delta errors are moderate. Boundary remains the weakest component, especially for the put case with $T=0.25$, $\sigma=0.4$, $r=0.01$, and $q=0.06$.

\begin{figure}[h]
\centering
\includegraphics[width=0.95\linewidth]{results/figures/stage11_integrated_workflow/integrated_price_rmse.png}
\caption{Price RMSE for each fresh PSOR audit regime. The positive-premium component preserves the payoff obstacle on these cases, but PSOR remains the reference label generator.}
\end{figure}

\begin{figure}[h]
\centering
\includegraphics[width=0.95\linewidth]{results/figures/stage11_integrated_workflow/integrated_delta_rmse.png}
\caption{Strict-mask Delta RMSE for each fresh PSOR audit regime. These errors evaluate the separate bounded Delta head and do not imply Gamma support.}
\end{figure}

% Auto-generated by experiments/19_report_visualization_enrichment.py
\section*{Visual Evidence Gallery}
\begin{figure}[!htbp]
\centering
\includegraphics[width=0.64\linewidth]{results/figures/report_enrichment/stage11_integrated_workflow_modular_workflow.pdf}
\caption{Separate price, boundary, and Delta heads support scenario screening only when paired with fresh CN/PSOR audit cases.}
\end{figure}

\begin{figure}[!htbp]
\centering
\begin{subfigure}[t]{0.48\linewidth}
\centering
\includegraphics[width=\linewidth]{results/figures/report_enrichment/stage11_integrated_workflow_integrated_price.pdf}
\caption{Per regime price error is retained to avoid hiding weak integrated cases.}
\end{subfigure}
\hfill
\begin{subfigure}[t]{0.48\linewidth}
\centering
\includegraphics[width=\linewidth]{results/figures/report_enrichment/stage11_integrated_workflow_integrated_boundary.pdf}
\caption{Boundary error remains a limitation in the integrated workflow evidence.}
\end{subfigure}
\end{figure}

\begin{figure}[!htbp]
\centering
\includegraphics[width=0.82\linewidth]{results/figures/report_enrichment/stage11_integrated_workflow_integrated_speed.pdf}
\caption{Speedups are reported alongside PSOR audit rather than as solver replacement evidence.}
\end{figure}

\FloatBarrier


\subsection{Failure Modes and Limitations}
The workflow is evaluated on only six deterministic held-out regimes, all with
short maturity, so it is not a deployment or universal-generalization study.
Boundary estimates remain materially weaker than price and Delta estimates.
PSOR remains required for audit, new regimes, and final label generation. This
stage does not report Gamma and does not merge the three heads into one model.

\subsection{Conclusion and Next Step}
The Stage 11 gate passes for modular scenario screening with PSOR audit. The next stage should move to claim audit and final report preparation, preserving the restricted language around the boundary head, the price model, and Delta diagnostics.

\end{document}
